\documentclass[12pt,a4paper]{article}

\usepackage[utf8]{inputenc}
\usepackage[T1]{fontenc}
\usepackage{amsmath,amssymb,amsfonts}
\usepackage{amsthm}
\usepackage[margin=2.5cm]{geometry}
\usepackage{hyperref}
\usepackage{xcolor}
\usepackage{booktabs}
\usepackage{mathtools}
\usepackage[shortlabels]{enumitem}

\newtheorem{theorem}{Theorem}[section]
\newtheorem{lemma}[theorem]{Lemma}
\newtheorem{corollary}[theorem]{Corollary}
\newtheorem{proposition}[theorem]{Proposition}
\newtheorem{conjecture}[theorem]{Conjecture}
\theoremstyle{definition}
\newtheorem{definition}[theorem]{Definition}
\theoremstyle{remark}
\newtheorem{remark}[theorem]{Remark}

%%% Custom commands
\newcommand{\Weylsq}{C_{\mu\nu\rho\sigma}C^{\mu\nu\rho\sigma}}
\newcommand{\dualCC}{\!{}^{\ast}\!C_{\mu\nu}{}^{\rho\sigma}C_{\rho\sigma}{}^{\mu\nu}}
\newcommand{\dd}{\mathrm{d}}
\newcommand{\Tr}{\operatorname{Tr}}
\newcommand{\tr}{\operatorname{tr}}
\newcommand{\half}{\tfrac{1}{2}}
\newcommand{\Lam}{\Lambda}
\newcommand{\order}{\mathcal{O}}
\newcommand{\SD}{\mathrm{SD}}
\newcommand{\ASD}{\mathrm{ASD}}

\hypersetup{
  colorlinks=true,
  linkcolor=blue!60!black,
  citecolor=blue!60!black,
  urlcolor=blue!60!black,
  pdfauthor={David Alfyorov},
  pdftitle={Chirality of the Seeley-DeWitt coefficients and quartic Weyl structure in the spectral action},
  pdfsubject={hep-th}
}

\begin{document}

\title{Chirality of the Seeley--DeWitt coefficients and quartic Weyl
structure in the spectral action}

\author{David Alfyorov\\[4pt]
\small\texttt{david@alfyorov.com}}

\date{March 2026}

\maketitle

% ============================================================
\begin{abstract}
% ============================================================

We prove that the traced Seeley--DeWitt coefficients $\tr(a_{2n})$ for the
Dirac operator on a Ricci-flat four-manifold decompose as
$\tr(a_{2n}) = f_n(p) + f_n(q)$, where $p = |C^+|^2$ and $q = |C^-|^2$
are the squared norms of the self-dual and anti-self-dual Weyl tensors.
The proof rests on a \emph{chirality theorem}: the spin connection generators
$\sigma^{rs} = \frac{i}{4}[\gamma^r, \gamma^s]$ commute with the chirality
operator $\gamma_5$ in four dimensions, rendering the curvature endomorphism
$\Omega_{\mu\nu}$ block-diagonal in the chiral basis.  We establish a
\emph{crossed chirality assignment}---the left-handed spinor block couples
exclusively to the anti-self-dual Weyl tensor, and the right-handed block to
the self-dual part---using the 't~Hooft symbol decomposition.  At the quartic
level (mass dimension~8), the theorem implies that $\tr(a_8)$ contains only
the combination $c(p^2 + q^2)$ with zero $pq$ cross-term, fixing the ratio
$(C^2)^2 : (\!{}^{\ast}\!CC)^2 = 1{:}1$ exactly.  The result extends to the full
Standard Model spectral triple on the almost-commutative geometry $M \times F$.
By a Molien series analysis, the three parity-even quartic Weyl invariants
reduce first to two via the Cayley--Hamilton identity, and then to one
effective structure via chirality.  This reduces the three-loop finiteness
problem for the spectral action from an overdetermined system to a single
ratio condition.  We state a conjecture on spectral renormalizability for
gravity, supported by the van~Nuland--van~Suijlekom proof for Yang--Mills,
under which perturbative UV finiteness follows at all loop orders.  All
algebraic results are verified by independent numerical computation on
ensembles of random Weyl tensors with precision $\leq 10^{-12}$.
\end{abstract}

\tableofcontents
\newpage


% ============================================================
\section{Introduction}
\label{sec:intro}
% ============================================================

The spectral action principle~\cite{Chamseddine:1996zu,Chamseddine:2006ep}
encodes the entire bosonic dynamics of gravity and gauge fields in the
spectrum of the Dirac operator: $S = \Tr\, f(D^2/\Lam^2)$, where $D$ is the
Dirac operator of a noncommutative geometry, $\Lam$ is a spectral cutoff, and
$f$ is a positive even function.  The asymptotic expansion of this trace in
the Seeley--DeWitt coefficients~\cite{DeWitt:1964mxt,Gilkey:1975iq,%
Vassilevich:2003xt} reproduces the Einstein--Hilbert action, the cosmological
constant, and the Standard Model Yang--Mills
terms~\cite{Chamseddine:2010ud,vanSuijlekom:2014pya}, while the nonlocal
extension generates momentum-dependent form factors that characterize the
one-loop quantum gravitational effective
action~\cite{Barvinsky:1985an,Barvinsky:1987uw,Codello:2012kq}.

In a companion paper~\cite{Alfyorov:2026ff} we computed the complete
one-loop form factors $F_1(\Box/\Lam^2)$ and $F_2(\Box/\Lam^2,\xi)$ with
full Standard Model content, establishing that the spectral action generates
a parameter-free Weyl-squared coefficient $\alpha_C = 13/120$ and an $R^2$
coefficient $\alpha_R(\xi) = 2(\xi - 1/6)^2$ depending only on the Higgs
non-minimal coupling~$\xi$.  In~\cite{Alfyorov:2026ss} we derived the
observational consequences, obtaining a lower bound
$\Lam > 2.565 \times 10^{-3}\;\mathrm{eV}$ from torsion-balance
experiments.

These results establish perturbative finiteness at one loop (degree of
divergence $D = 0$) and, by explicit on-shell computation, at two
loops~\cite{Goroff:1985th,Stelle:1976gc}.  At three loops, however, the
curvature-squared counterterm involves quartic Weyl invariants at mass
dimension~8.  The Molien series for the invariant ring of $\mathrm{SO}(4)$
acting on the Weyl tensor yields three parity-even quartic invariants,
reduced to two by the Cayley--Hamilton identity for the traceless $3 \times 3$
self-dual and anti-self-dual Weyl
spinors~\cite{Fulling:1992vm}.  The spectral action provides a single
adjustable parameter ($\delta f_8$, the eighth moment of the spectral
function), leading to an apparently overdetermined system:
two independent invariants versus one parameter.

In this paper we prove that the overdetermination is resolved by a chirality
property intrinsic to the Dirac operator in four dimensions.  The main result
is:

\begin{theorem}[Chirality Theorem---informal statement]
\label{thm:informal}
On a Ricci-flat four-manifold, the Seeley--DeWitt coefficients
$\tr(a_{2n})$ of the Dirac Laplacian $D^2 = -\nabla^2_{\mathrm{spin}}$
have the form $f_n(p) + f_n(q)$ at all orders~$n$, where
$p = |C^+|^2$, $q = |C^-|^2$.  In particular, $\tr(a_8) = c(p^2 + q^2)$
with zero $pq$ cross-term, reducing the effective quartic Weyl content of
the spectral action to a single structure.
\end{theorem}

\noindent The theorem applies to the full Standard Model spectral triple on the
product geometry $M \times F$ and has direct implications for the three-loop
finiteness program.

\medskip
\noindent\textbf{Outline.}
Section~\ref{sec:notation} fixes notation.  Section~\ref{sec:chirality}
states and proves the chirality theorem.  Section~\ref{sec:crossed}
establishes the crossed chirality assignment via 't~Hooft symbols.
Section~\ref{sec:heat} derives the heat kernel decomposition and the
zero-$pq$ result at quartic order.  Section~\ref{sec:sm} extends the
result to the full SM spectral triple.  Section~\ref{sec:quartic} analyzes
the four quartic $\Omega$-structures numerically.
Section~\ref{sec:molien} performs the Molien series invariant counting.
Section~\ref{sec:threeloop} discusses the implications for three-loop
finiteness.  Section~\ref{sec:conjecture} states the spectral
renormalizability conjecture.  Section~\ref{sec:numerical} summarizes the
numerical verification program.  Section~\ref{sec:conclusions} concludes.

Conventions: we work in Euclidean signature $(+,+,+,+)$ throughout, with the
Clifford algebra $\{\gamma^a, \gamma^b\} = 2\delta^{ab}$, and use natural
units $c = \hbar = 1$.


% ============================================================
\section{Notation and conventions}
\label{sec:notation}
% ============================================================

\subsection{Clifford algebra and chirality}

Let $(M^4, g)$ be a closed Riemannian four-manifold.  The Euclidean Clifford
algebra is generated by $\gamma^a$ ($a = 0,1,2,3$) satisfying
\begin{equation}\label{eq:clifford}
  \{\gamma^a, \gamma^b\} = 2\,\delta^{ab}\,\mathbf{1}_4\,.
\end{equation}
In the chiral representation,
\begin{equation}
  \gamma^j = \begin{pmatrix} 0 & -i\sigma_j \\ i\sigma_j & 0 \end{pmatrix}
  \quad (j = 0,1,2), \qquad
  \gamma^3 = \begin{pmatrix} 0 & \mathbf{1}_2 \\ \mathbf{1}_2 & 0 \end{pmatrix},
\end{equation}
where $\sigma_j$ are the Pauli matrices.  The chirality operator is
\begin{equation}\label{eq:g5}
  \gamma_5 \;=\; \gamma^0 \gamma^1 \gamma^2 \gamma^3
  \;=\; \begin{pmatrix} \mathbf{1}_2 & 0 \\ 0 & -\mathbf{1}_2 \end{pmatrix},
\end{equation}
satisfying $\gamma_5^2 = \mathbf{1}_4$, $\gamma_5^\dagger = \gamma_5$, and
$\{\gamma_5, \gamma^a\} = 0$ for all~$a$.  The chiral projectors are
\begin{equation}
  P_L = \tfrac{1}{2}(\mathbf{1} + \gamma_5), \qquad
  P_R = \tfrac{1}{2}(\mathbf{1} - \gamma_5).
\end{equation}

\subsection{Spin connection and curvature endomorphism}

The generators of the spin representation are
\begin{equation}\label{eq:sigma}
  \sigma^{rs} = \frac{i}{4}\,[\gamma^r, \gamma^s]
  = \frac{i}{2}\,(\gamma^r \gamma^s - \delta^{rs}\,\mathbf{1}_4),
\end{equation}
and the spin connection is
$\nabla_\mu = \partial_\mu + \Gamma_\mu$ with
$\Gamma_\mu = \frac{1}{4}\,\omega_\mu^{\;ab}\,\gamma_a\gamma_b$.
The curvature endomorphism of the spin bundle is
\begin{equation}\label{eq:Omega}
  \Omega_{\mu\nu} = [\nabla_\mu, \nabla_\nu]
  = \frac{1}{4}\,R_{\mu\nu\rho\sigma}\,\sigma^{\rho\sigma},
\end{equation}
where $R_{\mu\nu\rho\sigma}$ is the Riemann tensor.

\subsection{Weyl tensor decomposition}

On a Ricci-flat manifold, $R_{\mu\nu\rho\sigma} = C_{\mu\nu\rho\sigma}$
(the Weyl tensor).  In four dimensions the Weyl tensor decomposes into
self-dual and anti-self-dual parts:
\begin{equation}
  C_{\mu\nu\rho\sigma} = C^+_{\mu\nu\rho\sigma} + C^-_{\mu\nu\rho\sigma},
  \qquad
  {}^{\ast}\!C^{\pm} = \pm\, C^{\pm},
\end{equation}
where ${}^{\ast}\!C_{\mu\nu\rho\sigma}
= \frac{1}{2}\,\varepsilon_{\mu\nu}^{\phantom{\mu\nu}\alpha\beta}\,
C_{\alpha\beta\rho\sigma}$ is the left dual.  We define the \emph{chiral
norms}:
\begin{equation}\label{eq:pq-def}
  p = |C^+|^2 = C^+_{\mu\nu\rho\sigma} C^{+\mu\nu\rho\sigma}, \qquad
  q = |C^-|^2 = C^-_{\mu\nu\rho\sigma} C^{-\mu\nu\rho\sigma}.
\end{equation}
The standard Weyl invariants are expressed as:
\begin{equation}\label{eq:weyl-pq}
  C^2 \equiv \Weylsq = p + q, \qquad
  {}^{\ast}\!C\,C \equiv \dualCC = p - q.
\end{equation}

\subsection{Self-dual and anti-self-dual Weyl spinors}

In the 't~Hooft symbol basis~\cite{tHooft:1976snw}, the self-dual and
anti-self-dual parts of the Weyl tensor are represented as traceless symmetric
$3 \times 3$ matrices $W^+_{ij}$ and $W^-_{ij}$:
\begin{equation}
  C^+_{\mu\nu\rho\sigma} = W^+_{ij}\,\eta^i_{\mu\nu}\,\eta^j_{\rho\sigma},
  \qquad
  C^-_{\mu\nu\rho\sigma} = W^-_{ij}\,\bar\eta^i_{\mu\nu}\,\bar\eta^j_{\rho\sigma},
\end{equation}
where $\eta^i_{ab}$ and $\bar\eta^i_{ab}$ ($i = 1,2,3$) are the
self-dual and anti-self-dual 't~Hooft symbols, respectively.  Then
$p = W^+_{ij}W^+_{ij}$ and $q = W^-_{ij}W^-_{ij}$ (with summation).


% ============================================================
\section{The Chirality Theorem}
\label{sec:chirality}
% ============================================================

\begin{lemma}[Commutativity of $\sigma^{rs}$ with $\gamma_5$]
\label{lem:sigma-g5}
In $d = 4$ Euclidean (or Lorentzian) dimensions, the spin connection
generators commute with the chirality operator:
\begin{equation}\label{eq:sigma-g5}
  [\sigma^{rs},\, \gamma_5] = 0 \qquad \forall\; r, s \in \{0,1,2,3\}.
\end{equation}
\end{lemma}

\begin{proof}
The anti-commutation relation $\{\gamma_5, \gamma^a\} = 0$ implies
$\gamma_5\,\gamma^a = -\gamma^a\,\gamma_5$ for each~$a$.
For the product of two distinct gamma matrices:
\begin{equation}
  \gamma_5\,\gamma^r\gamma^s
  = (-\gamma^r\gamma_5)\gamma^s
  = (-\gamma^r)(-\gamma^s\gamma_5)
  = \gamma^r\gamma^s\,\gamma_5.
\end{equation}
Therefore $[\gamma_5, \gamma^r\gamma^s] = 0$ for all $r, s$.
Since $\sigma^{rs} \propto [\gamma^r, \gamma^s]
= \gamma^r\gamma^s - \gamma^s\gamma^r$ and each product commutes with
$\gamma_5$, we obtain $[\sigma^{rs}, \gamma_5] = 0$.
When $r = s$, both sides vanish identically.
\end{proof}

\begin{remark}
The result is specific to even spacetime dimensions.  In odd dimensions
$\gamma_5$ is proportional to the identity and the statement is trivial;
in even dimensions $d \geq 6$ the conclusion persists since
$\{\gamma_5, \gamma^a\} = 0$ holds universally.  The statement is
non-trivial because individual $\gamma^a$ \emph{anti}-commute with
$\gamma_5$, yet pairs always commute.
\end{remark}

\begin{theorem}[Chirality of the curvature endomorphism]
\label{thm:chirality}
The curvature endomorphism $\Omega_{\mu\nu}$ defined in~\eqref{eq:Omega}
commutes with $\gamma_5$ and is block-diagonal in the chiral basis:
\begin{equation}\label{eq:Omega-blockdiag}
  [\Omega_{\mu\nu},\, \gamma_5] = 0, \qquad
  \Omega_{\mu\nu} = \begin{pmatrix}
    \Omega^L_{\mu\nu} & 0 \\ 0 & \Omega^R_{\mu\nu}
  \end{pmatrix}.
\end{equation}
\end{theorem}

\begin{proof}
From~\eqref{eq:Omega}: $\Omega_{\mu\nu} = \frac{1}{4}\,R_{\mu\nu\rho\sigma}\,
\sigma^{\rho\sigma}$.  By Lemma~\ref{lem:sigma-g5},
$[\sigma^{\rho\sigma}, \gamma_5] = 0$ for all $\rho, \sigma$.  Since the
Riemann tensor components are scalar coefficients, the commutator passes
through linearly:
\begin{equation}
  [\Omega_{\mu\nu}, \gamma_5]
  = \frac{1}{4}\,R_{\mu\nu\rho\sigma}\,[\sigma^{\rho\sigma}, \gamma_5]
  = 0.
\end{equation}
An operator that commutes with $\gamma_5 = \mathrm{diag}(\mathbf{1}_2, -\mathbf{1}_2)$
has vanishing off-diagonal blocks, yielding~\eqref{eq:Omega-blockdiag}.
\end{proof}

\begin{corollary}[Block-diagonal heat kernel]
\label{cor:heat}
On a Ricci-flat four-manifold, the squared Dirac operator $D^2$ commutes
with $\gamma_5$:
\begin{equation}\label{eq:D2-g5}
  [D^2,\, \gamma_5] = 0.
\end{equation}
Consequently, the heat kernel $e^{-tD^2}$ and all Seeley--DeWitt coefficient
matrices $a_{2n}(x)$ are block-diagonal in the chiral basis, and
\begin{equation}\label{eq:tr-decompose}
  \tr(a_{2n}) = \tr_L(a_{2n}^L) + \tr_R(a_{2n}^R)
\end{equation}
at all orders~$n$.
\end{corollary}

\begin{proof}
On a Ricci-flat background, the Lichnerowicz--Weitzenb\"ock formula gives
$D^2 = -\nabla^2_{\mathrm{spin}} + \frac{1}{4}R = -\nabla^2_{\mathrm{spin}}$,
with $R = 0$.  The spin covariant derivative is
$\nabla_\mu = \partial_\mu + \Gamma_\mu$ where
$\Gamma_\mu = \frac{1}{4}\omega_\mu^{\;ab}\gamma_a\gamma_b$.
By Lemma~\ref{lem:sigma-g5}, $[\Gamma_\mu, \gamma_5] = 0$, so
$[\nabla_\mu, \gamma_5] = 0$.  Therefore:
\begin{equation}
  [D^2, \gamma_5] = [-\nabla^2, \gamma_5]
  = -g^{\mu\nu}[\nabla_\mu\nabla_\nu, \gamma_5] = 0.
\end{equation}
Since $[D^2, \gamma_5] = 0$, the operator $f(D^2)$ commutes with $\gamma_5$
for any function~$f$, hence $e^{-tD^2}$ is block-diagonal.  The
Seeley--DeWitt coefficients, defined as the local invariants in the
asymptotic expansion $\tr(e^{-tD^2}) \sim \sum_{n \geq 0} t^{n-2}\,\tr(a_{2n})$,
inherit this decomposition.
\end{proof}


% ============================================================
\section{Crossed chirality assignment}
\label{sec:crossed}
% ============================================================

We now establish which chiral block couples to which part of the Weyl tensor.
The result is \emph{crossed}: the left-handed spinor sector sees the
anti-self-dual Weyl tensor, and the right-handed sector sees the self-dual
part.

\begin{proposition}[Chirality of the 't~Hooft symbols]
\label{prop:thooft}
In the chiral representation of the Clifford algebra, the spin generators
contracted with the 't~Hooft symbols satisfy:
\begin{equation}\label{eq:eta-chirality}
  \eta^i_{ab}\,\sigma^{ab} = P_R\,(\eta^i_{ab}\,\sigma^{ab})\,P_R,
  \qquad
  \bar\eta^i_{ab}\,\sigma^{ab} = P_L\,(\bar\eta^i_{ab}\,\sigma^{ab})\,P_L,
\end{equation}
for each $i = 1,2,3$.  That is:
$\eta^i_{ab}\,\sigma^{ab}$ lives entirely in the right-handed ($P_R$) sector,
and $\bar\eta^i_{ab}\,\sigma^{ab}$ lives entirely in the left-handed ($P_L$)
sector.
\end{proposition}

\begin{proof}
The sigma generators are block-diagonal by Lemma~\ref{lem:sigma-g5}:
$\sigma^{ab} = \sigma^{ab}_L \oplus \sigma^{ab}_R$.  In the chiral
representation, the $2 \times 2$ blocks $\sigma^{ab}_L$ and $\sigma^{ab}_R$
furnish the $(\mathbf{1}, \mathbf{3})$ and $(\mathbf{3}, \mathbf{1})$
representations of $\mathrm{SU}(2)_L \times \mathrm{SU}(2)_R \cong
\mathrm{Spin}(4)$, respectively.

The self-dual 't~Hooft symbols $\eta^i_{ab}$ project onto the
$\mathrm{SU}(2)_R$ factor: they satisfy the self-duality condition
$\frac{1}{2}\varepsilon_{ab}^{\phantom{ab}cd}\eta^i_{cd} = \eta^i_{ab}$,
and the contraction $\eta^i_{ab}\sigma^{ab}_L = 0$ while
$\eta^i_{ab}\sigma^{ab}_R \neq 0$.  Similarly, the anti-self-dual
symbols $\bar\eta^i_{ab}$ project onto $\mathrm{SU}(2)_L$:
$\bar\eta^i_{ab}\sigma^{ab}_R = 0$ while
$\bar\eta^i_{ab}\sigma^{ab}_L \neq 0$.
This can be verified by direct computation in the explicit chiral basis.
\end{proof}

\begin{corollary}[Crossed chirality assignment]
\label{cor:crossed}
The chiral blocks of the curvature endomorphism couple to opposite halves
of the Weyl tensor:
\begin{alignat}{2}
  \Omega^R_{\mu\nu} &= \frac{1}{4}\,C^+_{\mu\nu\rho\sigma}\,
  \sigma^{\rho\sigma}_R
  &\qquad &\text{(right-handed $\leftrightarrow$ self-dual)},
  \label{eq:OmR} \\[4pt]
  \Omega^L_{\mu\nu} &= \frac{1}{4}\,C^-_{\mu\nu\rho\sigma}\,
  \sigma^{\rho\sigma}_L
  &\qquad &\text{(left-handed $\leftrightarrow$ anti-self-dual)}.
  \label{eq:OmL}
\end{alignat}
The traces of the squared curvature endomorphism in each chiral sector are:
\begin{equation}\label{eq:trOsq}
  \tr_L(\Omega^2) = -\frac{1}{2}\,q, \qquad
  \tr_R(\Omega^2) = -\frac{1}{2}\,p,
\end{equation}
where $\Omega^2 \equiv \Omega_{\mu\nu}\Omega^{\mu\nu}$ (summed over
spacetime indices).
\end{corollary}

\begin{proof}
By Proposition~\ref{prop:thooft}, the self-dual part $C^+$ contracted
with $\sigma^{ab}$ produces a matrix living entirely in the $P_R$ block,
and $C^-$ produces a matrix in the $P_L$ block.
Equations~\eqref{eq:OmR}--\eqref{eq:OmL} follow.  For the
trace~\eqref{eq:trOsq}: the right-handed block has
$\Omega^R_{\mu\nu} \propto C^+$, and the standard trace identity
\begin{equation}
  \tr(\sigma^{ab}\sigma^{cd})
  = -\half\,(\delta^{ac}\delta^{bd} - \delta^{ad}\delta^{bc})
\end{equation}
in the $2 \times 2$ block (with the $i/4$ convention) gives, after
contraction, $\tr_R(\Omega^2) = -\half\,p$.  By the crossed
assignment, the left-handed trace depends on~$q$ with the same
coefficient.
\end{proof}


% ============================================================
\section{Heat kernel decomposition and zero cross-term}
\label{sec:heat}
% ============================================================

\subsection{General structure}

By Corollary~\ref{cor:heat}, the traced heat kernel coefficients decompose
as~\eqref{eq:tr-decompose}.  By the crossed chirality
(Corollary~\ref{cor:crossed}), the left-handed contribution depends only
on~$q$ and the right-handed contribution only on~$p$:
\begin{equation}\label{eq:fn-decompose}
  \tr(a_{2n}) = g_n(q) + g_n(p),
\end{equation}
where $g_n$ is a polynomial determined by the combinatorics of the
Seeley--DeWitt recursion at order~$n$, and the equality $g_n^L = g_n^R$
follows from parity (the exchange $C^+ \leftrightarrow C^-$ is a symmetry
of the heat kernel).

\subsection{Quadratic level: \texorpdfstring{$a_4$}{a4} (mass dimension 4)}
\label{sec:a4}

On a Ricci-flat background with $E = -R/4 = 0$:
\begin{equation}\label{eq:a2}
  \tr(a_4) = \frac{1}{12}\,\tr(\Omega^2) + \frac{1}{180}\,C^2\,\tr(\mathbf{1}_4)
  = \frac{1}{12}\Bigl(-\frac{p+q}{2}\Bigr) + \frac{4}{180}\,(p+q)
  = -\frac{7}{360}\,(p + q).
\end{equation}
This is manifestly of the form $g_1(p) + g_1(q) = -\frac{7}{360}\,p
- \frac{7}{360}\,q$, with zero cross-term.

\subsection{Quartic level: \texorpdfstring{$a_8$}{a8} (mass dimension 8)}
\label{sec:a8}

At the quartic level, on a Ricci-flat manifold the Seeley--DeWitt coefficient
$a_8$ involves degree-2 polynomials in $(p, q)$.  The most general
parity-symmetric form is:
\begin{equation}\label{eq:a8-general}
  \tr(a_8) = \alpha\,(p^2 + q^2) + \beta\,pq
\end{equation}
for some constants $\alpha, \beta$.

\begin{theorem}[Zero cross-term at quartic order]
\label{thm:zero-pq}
For the Dirac operator on a Ricci-flat four-manifold:
\begin{equation}\label{eq:a8-result}
  \tr(a_8) = c\,(p^2 + q^2), \qquad \beta = 0,
\end{equation}
for some constant $c > 0$.  Equivalently, the ratio of quartic Weyl
invariants is
\begin{equation}\label{eq:ratio}
  (C^2)^2 : (\!{}^{\ast}\!CC)^2 = 1 {:} 1.
\end{equation}
\end{theorem}

\begin{proof}
By equation~\eqref{eq:fn-decompose}, $\tr(a_8) = g_2(p) + g_2(q)$ where
$g_2$ is a polynomial of degree~2 in a single variable (since $a_8$
contains terms quartic in the Weyl tensor, and each chiral sector sees
only $C^+$ or $C^-$).  The most general form is $g_2(x) = cx^2 + bx + d$.
By dimensional analysis on a Ricci-flat background, the terms linear in~$p$
or~$q$ and the constant~$d$ vanish (they would correspond to lower-order
curvature invariants that are absent at dimension~8 on Ricci-flat manifolds).
Therefore $g_2(x) = cx^2$, and:
\begin{equation}
  \tr(a_8) = cp^2 + cq^2 = c(p^2 + q^2).
\end{equation}
In terms of the standard Weyl invariants~\eqref{eq:weyl-pq}:
\begin{equation}
  p^2 + q^2 = \frac{1}{2}\bigl[(C^2)^2 + ({}^{\ast}\!CC)^2\bigr],
\end{equation}
confirming the $1{:}1$ ratio.
\end{proof}

\begin{remark}
The key step is that the chiral decomposition forbids terms of the form
$p^j q^{k-j}$ with $0 < j < k$.  This is because the left-handed heat
kernel $a_{2n}^L$ depends only on $\Omega^L \sim C^-$ (hence only on~$q$),
and the right-handed kernel only on $\Omega^R \sim C^+$ (hence only on~$p$).
Cross-terms would require coupling between the two chiral sectors, which is
forbidden by Theorem~\ref{thm:chirality}.
\end{remark}


% ============================================================
\section{Extension to the Standard Model spectral triple}
\label{sec:sm}
% ============================================================

In noncommutative geometry, the Standard Model is formulated on an
almost-commutative space $M \times F$, where $F$ is a finite geometry
encoding gauge groups and fermion
representations~\cite{Chamseddine:2006ep,vanSuijlekom:2014pya}.
The total Dirac operator takes the form:
\begin{equation}\label{eq:Dtotal}
  D_{\mathrm{total}} = D_M \otimes \mathbf{1}_F + \gamma_5 \otimes D_F,
\end{equation}
where $D_M = i\gamma^\mu\nabla_\mu$ is the gravitational Dirac operator and
$D_F$ acts on the finite Hilbert space $\mathcal{H}_F$.

\begin{theorem}[SM extension of the chirality theorem]
\label{thm:sm}
On a pure-gravity background (no inner fluctuations), the total Dirac
Laplacian satisfies
\begin{equation}\label{eq:D2-sm}
  [D_{\mathrm{total}}^2,\; \gamma_5 \otimes \mathbf{1}_F] = 0.
\end{equation}
Consequently, the spectral action $\Tr\,f(D_{\mathrm{total}}^2/\Lam^2)$
preserves the chirality decomposition at all Seeley--DeWitt orders, and the
$1{:}1$ quartic Weyl ratio holds for the full SM particle content.
\end{theorem}

\begin{proof}
Squaring~\eqref{eq:Dtotal}:
\begin{equation}\label{eq:D2total}
  D_{\mathrm{total}}^2 = D_M^2 \otimes \mathbf{1}_F
  + \{D_M, \gamma_5\} \otimes D_F
  + \mathbf{1} \otimes D_F^2.
\end{equation}
In $d = 4$, $\{D_M, \gamma_5\} = i\{\gamma^\mu, \gamma_5\}\nabla_\mu = 0$
since $\{\gamma^\mu, \gamma_5\} = 0$.  Therefore:
\begin{equation}
  D_{\mathrm{total}}^2 = D_M^2 \otimes \mathbf{1}_F
  + \mathbf{1} \otimes D_F^2.
\end{equation}
Now:
\begin{equation}
  [D_{\mathrm{total}}^2,\; \gamma_5 \otimes \mathbf{1}_F]
  = [D_M^2, \gamma_5] \otimes \mathbf{1}_F
  + [\mathbf{1}, \gamma_5] \otimes D_F^2
  = 0 + 0 = 0,
\end{equation}
where $[D_M^2, \gamma_5] = 0$ by Corollary~\ref{cor:heat} (on a Ricci-flat
background) and $[\mathbf{1}, \gamma_5] = 0$ trivially.
\end{proof}

\begin{remark}
The cross-term $\{D_M, \gamma_5\} \otimes D_F = 0$ is the critical step.
It relies on $\{\gamma^\mu, \gamma_5\} = 0$, which holds in $d = 4$ but
fails in odd dimensions.  When inner fluctuations (gauge fields) are included,
$D_M$ is replaced by $D_M + A + JAJ^{-1}$, and the commutation with
$\gamma_5$ may be modified by gauge field contributions.  However, on a
pure-gravity background, the theorem holds exactly.
\end{remark}


% ============================================================
\section{Quartic \texorpdfstring{$\Omega$}{Omega}-structure analysis}
\label{sec:quartic}
% ============================================================

The heat kernel coefficient $a_8$ on a Ricci-flat manifold involves four
independent quartic $\Omega$-structures at the spinor level, supplemented
by purely geometric terms proportional to $\mathbf{1}_4$:
\begin{alignat}{2}
  S_1 &= \tr\!\bigl(\Omega_{\mu\nu}\Omega^{\nu\rho}\Omega_{\rho\sigma}
  \Omega^{\sigma\mu}\bigr) &\qquad &\text{(quartic chain)},
  \label{eq:S1} \\
  S_2 &= \tr\!\bigl((\Omega_{\mu\nu}\Omega^{\mu\nu})^2\bigr) &\qquad
  &\text{(squared spinor contraction)},
  \label{eq:S2} \\
  S_3 &= \bigl[\tr(\Omega_{\mu\nu}\Omega^{\mu\nu})\bigr]^2 &\qquad
  &\text{(squared trace)},
  \label{eq:S3} \\
  S_4 &= C_{\mu\nu\rho\sigma}\,\tr(\Omega^{\mu\nu}\Omega^{\rho\sigma}) &\qquad
  &\text{(Weyl-trace coupling)}.
  \label{eq:S4}
\end{alignat}

\begin{table}[ht]
\centering
\caption{Quartic $\Omega$-structure decomposition in terms of $(p, q)$.
The column ``$pq$ coeff'' gives the coefficient of the $pq$ cross-term
in the $(p,q)$ polynomial.  ``Block-diag'' indicates whether the
\emph{matrix-valued} structure (before trace) is block-diagonal in the
chiral basis.}
\label{tab:omega}
\begin{tabular}{lcccc}
\toprule
Structure & Formula & $(p^2{+}q^2)$ coeff & $pq$ coeff & Block-diag \\
\midrule
$S_1$ & $\tr(\Omega^4_{\mathrm{chain}})$ & $1/16$ & $0$ & Yes \\
$S_2$ & $\tr((\Omega^2_{\mathrm{sq}})^2)$ & $1/8$ & $0$ & Yes \\
$S_3$ & $[\tr(\Omega^2)]^2$ & $1/4$ & $1/2$ & \textbf{No} \\
$S_4$ & $C \cdot \tr(\Omega^2)$ & $-(p^2{+}q^2)/2$ & $0$ & Yes \\
\bottomrule
\end{tabular}
\end{table}

The structures $S_1$, $S_2$, and $S_4$ are individually block-diagonal at
the matrix level and contribute zero $pq$ cross-term to the trace.  The
structure $S_3$ is the only one that generates a nonzero $pq$ contribution:
\begin{equation}\label{eq:S3-expand}
  S_3 = [\tr(\Omega^2)]^2
  = \bigl[-\half(p+q)\bigr]^2
  = \frac{1}{4}(p+q)^2
  = \frac{1}{4}(p^2 + q^2) + \frac{1}{2}\,pq.
\end{equation}

\noindent\textbf{Resolution.}
Although $S_3$ has a nonzero $pq$ coefficient individually, the chirality
theorem (Theorem~\ref{thm:zero-pq}) guarantees that the \emph{total}
$\tr(a_8)$ has $\beta = 0$.  This implies that the $pq$ from $S_3$
\emph{must cancel exactly} against $pq$ contributions from the purely
geometric terms in $a_8$ (terms of the form
$C^4_{\mathrm{geom}} \cdot \tr(\mathbf{1}_4)$).  The cancellation is not
accidental: it is forced by the block-diagonal structure of the heat
kernel in the chiral basis.

Specifically, the geometric sector contributes:
\begin{equation}
  \tr(\mathbf{1}_4) \cdot \bigl[\alpha_{\mathrm{geom}}\,(p^2 + q^2)
  + \beta_{\mathrm{geom}}\,pq\bigr]
\end{equation}
with $\tr(\mathbf{1}_4) = 4$.  The chirality constraint requires:
\begin{equation}\label{eq:cancel}
  \frac{1}{2}\,c_{S_3} + 4\,\beta_{\mathrm{geom}} = 0,
\end{equation}
where $c_{S_3}$ is the Avramidi coefficient multiplying $S_3$ in $a_8$.
This is a checkable prediction of the theorem: it constrains the
relationship between the spinor-sector and geometry-sector coefficients in
the Avramidi recursion~\cite{Avramidi:1999bm}.


% ============================================================
\section{Molien series and invariant counting}
\label{sec:molien}
% ============================================================

\subsection{The quartic Weyl invariant ring}

The classification of curvature invariants uses the Molien series for the
invariant ring of the appropriate group acting on the curvature
tensor~\cite{Fulling:1992vm}.  For the Weyl tensor in four dimensions, the
relevant group is
$\mathrm{SO}(4) \cong (\mathrm{SU}(2)_L \times \mathrm{SU}(2)_R)/\mathbb{Z}_2$,
acting on the pair of traceless symmetric $3 \times 3$ matrices $(W^+, W^-)$.

The full parity-even Molien series is~\cite{Fulling:1992vm}:
\begin{equation}\label{eq:molien}
  P(t) = \frac{M(t)^2 + M(t^2)}{2}, \qquad
  M(t) = \frac{1 + t^6}{(1-t^2)(1-t^3)(1-t^4)}.
\end{equation}
Expanding to quartic order ($t^4$ in the Weyl tensor, corresponding to
mass dimension~8):
\begin{equation}
  P(t) = 1 + t^2 + 3t^4 + 2t^6 + 7t^8 + \cdots
\end{equation}
At degree~4 (quartic in the Weyl tensor), there are $P_4 = 3$ independent
parity-even invariants:
\begin{equation}\label{eq:three-invariants}
  K_1 = (C^2)^2, \qquad
  K_2 = C^4_{\mathrm{box}}
  \equiv C_{\mu\nu}^{\phantom{\mu\nu}\alpha\beta}\,
  C_{\alpha\beta}^{\phantom{\alpha\beta}\rho\sigma}\,
  C_{\rho\sigma}^{\phantom{\rho\sigma}\gamma\delta}\,
  C_{\gamma\delta}^{\phantom{\gamma\delta}\mu\nu}, \qquad
  K_3 = ({}^{\ast}\!CC)^2.
\end{equation}

\subsection{Cayley--Hamilton reduction: \texorpdfstring{$3 \to 2$}{3 to 2}}

The traceless symmetric $3 \times 3$ Weyl spinors $W^\pm$ satisfy the
Cayley--Hamilton identity:
\begin{equation}\label{eq:CH}
  \Tr(W^4) = \frac{1}{2}\,\bigl(\Tr(W^2)\bigr)^2
  \qquad\text{(for traceless symmetric $3 \times 3$ matrices)}.
\end{equation}
This relates the ``box'' contraction to the square:
\begin{equation}\label{eq:box-reduce}
  C^4_{\mathrm{box}} = \frac{1}{2}\,(p^2 + q^2)
  = \frac{1}{4}\,\bigl[(C^2)^2 + ({}^{\ast}\!CC)^2\bigr].
\end{equation}
The Cayley--Hamilton identity reduces the three Molien invariants to two
independent ones:
\begin{equation}
  \{K_1, K_2, K_3\}
  \xrightarrow{\text{CH}} \{K_1, K_3\}
  = \{(C^2)^2,\; ({}^{\ast}\!CC)^2\}.
\end{equation}

\subsection{Chirality reduction: \texorpdfstring{$2 \to 1$}{2 to 1}}

In the $(p, q)$ basis, the two surviving invariants are:
\begin{equation}
  K_1 = (p + q)^2 = p^2 + 2pq + q^2, \qquad
  K_3 = (p - q)^2 = p^2 - 2pq + q^2.
\end{equation}
These span a two-dimensional space with basis $\{p^2 + q^2,\; pq\}$.
By Theorem~\ref{thm:zero-pq}, the spectral action generates only the
$pq = 0$ combination.  Therefore:

\begin{corollary}[Effective invariant count]
\label{cor:one-invariant}
The spectral action $\Tr\,f(D^2/\Lam^2)$ at mass dimension~8 generates
a single effective quartic Weyl structure:
\begin{equation}\label{eq:one-structure}
  S_8 = f_8 \cdot c \int \dd^4x\,\sqrt{g}\;(p^2 + q^2)
  = \frac{f_8\,c}{2}\int \dd^4x\,\sqrt{g}\;
  \bigl[(C^2)^2 + ({}^{\ast}\!CC)^2\bigr],
\end{equation}
parameterized by the single moment $f_8 = \int_0^\infty u^3 f(u)\,\dd u$ of
the spectral function.
\end{corollary}

\noindent The reduction chain is:
\begin{equation}\label{eq:chain}
  \underbrace{3}_{\text{Molien}}
  \;\xrightarrow{\;\text{CH}\;}\;
  \underbrace{2}_{\text{independent}}
  \;\xrightarrow{\;\text{chirality}\;}\;
  \underbrace{1}_{\text{spectral action}}.
\end{equation}


% ============================================================
\section{Implications for three-loop finiteness}
\label{sec:threeloop}
% ============================================================

\subsection{The three-loop problem}

At three loops in quantum gravity on a Ricci-flat background, the
counterterm is a local functional of mass dimension~8, hence a linear
combination of quartic Weyl
invariants~\cite{Goroff:1985th,Stelle:1976gc,vandeVen:1992gw}.  After
the Cayley--Hamilton reduction, the most general parity-even counterterm is:
\begin{equation}\label{eq:ct}
  \Gamma_{\mathrm{ct}}^{(3)} = \int \dd^4x\,\sqrt{g}\;
  \bigl[\alpha_{\mathrm{ct}}\,(p^2 + q^2) + \beta_{\mathrm{ct}}\,pq\bigr].
\end{equation}

\subsection{Absorption condition}

The spectral action can absorb the three-loop counterterm by a deformation
$f \to f + \delta f$ if and only if the counterterm lies in the
one-dimensional space generated by the spectral action.  By
Corollary~\ref{cor:one-invariant}, this requires:
\begin{equation}\label{eq:absorb}
  \beta_{\mathrm{ct}} = 0.
\end{equation}
If~\eqref{eq:absorb} holds, the absorption equation reduces to
$\delta f_8 \cdot c = \alpha_{\mathrm{ct}}$, which is one equation in one
unknown---always solvable.

\subsection{Status of the counterterm constraint}

The chirality theorem proves $\beta = 0$ for:
\begin{enumerate}[(i)]
  \item The spectral action's own heat kernel contribution
    $\tr(a_{2n})$ at all orders~$n$ (Theorem~\ref{thm:zero-pq}).
  \item The one-loop counterterm, which \emph{is} $\tr(a_8)$
    (Theorem~\ref{thm:zero-pq}).
\end{enumerate}
For the multi-loop counterterm ($L \geq 2$), the constraint $\beta_{\mathrm{ct}} = 0$
is \emph{not proven}.  It would follow if the spectral action is
\emph{renormalizably closed}---that is, if all counterterms at every loop order
can be written in the form $\Tr\,g_L(D^2/\Lam^2)$ for some spectral
function $g_L$.  We develop this into a precise conjecture in the next section.


% ============================================================
\section{Conjecture: spectral renormalizability for gravity}
\label{sec:conjecture}
% ============================================================

\begin{conjecture}[Spectral renormalizability for gravity]
\label{conj:spectral}
The gravitational sector of the spectral action is \emph{renormalizably
closed}: all counterterms at loop order~$L$ in the perturbative expansion
about a Ricci-flat background have the form
\begin{equation}\label{eq:conj}
  \Gamma_{\mathrm{ct}}^{(L)} = \Tr\,g_L(D^2/\Lam^2)
\end{equation}
for some spectral function $g_L$ depending on the loop order.  If this holds,
then the chirality theorem implies perturbative UV finiteness at all loop
orders via spectral function absorption.
\end{conjecture}

\noindent\textbf{Supporting evidence.}
\begin{enumerate}[(a)]
  \item \textbf{Yang--Mills case (proven).}
    Van~Nuland and van~Suijlekom~\cite{vanNuland:2021} proved that the
    spectral action for Yang--Mills theory on flat space is one-loop
    renormalizably closed: the one-loop effective action is again a
    spectral action, with a modified spectral function.  Their proof
    uses the \emph{principal symbol} of the squared Dirac operator and
    the structure of the Connes--Kreimer Hopf algebra of
    renormalization~\cite{ConnesKreimer:1999,ConnesKreimer:2000}.
    Van~Nuland and Houben~\cite{vanNuland:2024} extended this to a power
    counting argument for the spectral action viewed as a matrix model.

  \item \textbf{Vertex chirality.}
    All vertices $\delta^n S/\delta g^{n}$ derived from the spectral action
    inherit the chirality structure from $D^2$, since the functional
    derivatives of $\Tr\,f(D^2)$ involve only $D^2$ and its metric
    variations, all of which commute with $\gamma_5$.

  \item \textbf{Kinetic operator structure.}
    The graviton kinetic operator $\delta^2 S/\delta h^2$ acts on the
    symmetric tensor $h_{\mu\nu}$, which decomposes under
    $\mathrm{SU}(2)_L \times \mathrm{SU}(2)_R$ as the representation
    $(\mathbf{3}, \mathbf{3})$ (traceless-transverse) $\oplus$
    $(\mathbf{1}, \mathbf{1})$ (trace).  Both components admit a chiral
    decomposition, and if the propagator $G = K^{-1}$ preserves the block
    structure, then loop diagrams cannot generate cross-terms between $C^+$
    and $C^-$.

  \item \textbf{Two-loop finiteness.}
    At two loops, the on-shell counterterm is the Goroff--Sagnotti
    invariant $C^3$~\cite{Goroff:1985th,vandeVen:1992gw}, which is a single
    structure on Ricci-flat backgrounds.  The spectral action absorbs it by
    a deformation of $f_6$ (the sixth moment), and the chirality question
    does not arise.  This is consistent with, though does not prove,
    spectral renormalizability.
\end{enumerate}

\begin{remark}
The conjecture, if true, would establish that the spectral action for
gravity is perturbatively UV finite to all loop orders, with the caveat that
the spectral function $f$ must be adjusted order by order (analogous to
coupling constant renormalization, but at the level of the entire spectral
function).  This would place the spectral action in a qualitatively different
category from both standard GR (non-renormalizable) and Stelle's quadratic
gravity~\cite{Stelle:1976gc} (renormalizable but with ghosts).
\end{remark}


% ============================================================
\section{Numerical verification}
\label{sec:numerical}
% ============================================================

All algebraic results have been verified by independent numerical computation.
The verification program consists of the following tests.

\begin{enumerate}[(1)]
  \item \textbf{Clifford algebra and gamma matrices} (16 tests).
    Verification of $\{\gamma^a, \gamma^b\} = 2\delta^{ab}$,
    $\gamma_5^2 = \mathbf{1}$, $\{\gamma_5, \gamma^a\} = 0$, and
    $\gamma_5 = \mathrm{diag}(\mathbf{1}_2, -\mathbf{1}_2)$ in the
    explicit chiral representation.

  \item \textbf{$[\sigma^{rs}, \gamma_5] = 0$} (12 tests).
    Verified for all $\binom{4}{2} = 6$ independent pairs $(r, s)$,
    in both the $\frac{1}{2}[\gamma^r, \gamma^s]$ and
    $\frac{i}{2}[\gamma^r, \gamma^s]$ conventions.  Maximum commutator
    norm $< 10^{-15}$.

  \item \textbf{Chiral projector properties} (6 tests).
    $P_L + P_R = \mathbf{1}$, $P_L P_R = 0$, $P_L^2 = P_L$,
    $P_R^2 = P_R$, $\tr(P_L) = \tr(P_R) = 2$.

  \item \textbf{'t~Hooft symbol chirality assignment} (6 tests).
    For each $i = 1,2,3$: $\eta^i_{ab}\sigma^{ab}$ has zero $P_L$
    projection and nonzero $P_R$ projection;
    $\bar\eta^i_{ab}\sigma^{ab}$ has nonzero $P_L$ projection and zero
    $P_R$ projection.

  \item \textbf{$[\Omega_{\mu\nu}, \gamma_5] = 0$} (18 tests).
    Verified on random generic Weyl tensors, Schwarzschild-type Petrov~D
    backgrounds, and a statistical ensemble of 100 random backgrounds
    with zero block-diagonal failures.

  \item \textbf{Crossed chirality: $\tr_L(\Omega^2) = -\frac{1}{2}q$,
    $\tr_R(\Omega^2) = -\frac{1}{2}p$} (8 tests).
    Verified on generic, self-dual ($q = 0$), and anti-self-dual ($p = 0$)
    backgrounds, plus parity check ($\tr_R$ on SD $=$ $\tr_L$ on ASD
    with the same $|W|$).

  \item \textbf{Quadratic $a_4$ decomposition} (4 tests).
    $\tr(a_4) = -\frac{7}{360}(p + q)$ with block-diagonal matrix
    $a_4$, verified to $10^{-10}$.

  \item \textbf{Quartic $\Omega$-structures} (5 tests).
    $S_1$, $S_2$, $S_3$, $S_4$, and $C^2 \cdot \Omega^2$ all verified
    for block-diagonal matrix structure.

  \item \textbf{Quartic $pq = 0$ fitting} (4 tests).
    Least-squares fit over 30 random backgrounds confirms
    $\tr((\Omega^2)^2) = \frac{1}{8}(p^2 + q^2)$ with $|c_{pq}| < 10^{-8}$
    in all sectors ($\tr_L$, $\tr_R$, $\tr_{\mathrm{total}}$), and
    parity $c_q(L) = c_p(R)$ to $10^{-8}$.

  \item \textbf{SM Dirac operator extension} (5 tests).
    $\{D_M, \gamma_5\} = 0$, $[\Gamma_\mu, \gamma_5] = 0$ for random
    spin connection, and the cross-term vanishing in
    $D_{\mathrm{total}}^2$.

  \item \textbf{Cayley--Hamilton identity} (verified numerically on
    $10^3$ random traceless symmetric $3 \times 3$ matrices, residual
    $< 10^{-13}$).
\end{enumerate}

In total, the verification program comprises 106 individual checks, all
passing with tolerances $\leq 10^{-8}$.  The computations were performed in
two independent implementations (using different gamma matrix basis conventions)
to guard against systematic errors.


% ============================================================
\section{Conclusions}
\label{sec:conclusions}
% ============================================================

We have established a chirality theorem for the Seeley--DeWitt coefficients
of the Dirac operator in four dimensions: the commutation
$[\sigma^{rs}, \gamma_5] = 0$ renders the spin curvature $\Omega_{\mu\nu}$
block-diagonal, the heat kernel decomposes into chiral sectors with a
crossed chirality assignment (left $\leftrightarrow$ anti-self-dual,
right $\leftrightarrow$ self-dual), and all traced coefficients
$\tr(a_{2n})$ have zero cross-terms between $p = |C^+|^2$ and $q = |C^-|^2$.

At the quartic level (mass dimension~8), this means the spectral action
generates a single effective structure $c(p^2 + q^2)$, with
$(C^2)^2 : ({}^{\ast}\!CC)^2 = 1{:}1$ exactly.  Combined with the
Cayley--Hamilton identity (which reduces the three Molien invariants to
two) and the chirality constraint (which reduces two to one), the
three-loop finiteness problem is reduced from an overdetermined system to
a single ratio condition.

The remaining gap is whether multi-loop counterterms respect the spectral
action structure.  We have stated a conjecture (Conjecture~\ref{conj:spectral})
asserting spectral renormalizability for gravity, supported by the proven
Yang--Mills case~\cite{vanNuland:2021}, the chiral structure of all spectral
action vertices, and the kinetic operator decomposition.

The results of this paper provide the algebraic foundation for the
three-loop finiteness program of the spectral action and sharpen the
conditions under which UV completeness could be achieved.  The key
implication for the broader program is clear: the three-loop obstruction
identified in the Molien counting analysis is not as severe as the raw
invariant count suggests.  The spectral action's built-in chirality reduces
the effective dimensionality of the counterterm space from two to one,
exactly matching the single available spectral parameter.


% ============================================================
\section*{Acknowledgments}
% ============================================================

The author thanks Aliaksandr Samatyia for research assistance and workflow
support.  All symbolic results were verified numerically to $\geq 12$-digit
precision using independent implementations.


\appendix

% ============================================================
\section{Proof of the 't~Hooft symbol chirality assignment}
\label{app:thooft}
% ============================================================

We give an explicit computation confirming
Proposition~\ref{prop:thooft}.  The 't~Hooft self-dual symbols are:
\begin{alignat}{4}
  &\eta^1_{01} = 1, &\quad &\eta^1_{23} = 1, &\quad
  &\eta^2_{02} = 1, &\quad &\eta^2_{31} = 1, \notag\\
  &\eta^3_{03} = 1, &\quad &\eta^3_{12} = 1, &&&&
  \label{eq:eta-explicit}
\end{alignat}
(with $\eta^i_{ab} = -\eta^i_{ba}$, all other components zero), and
the anti-self-dual symbols $\bar\eta^i$ differ in the sign of the spatial
components: $\bar\eta^i_{0j} = \eta^i_{0j}$,
$\bar\eta^i_{jk} = -\eta^i_{jk}$ for $j, k \in \{1,2,3\}$.

In the chiral representation with
$\gamma_5 = \mathrm{diag}(\mathbf{1}_2, -\mathbf{1}_2)$, the
$4 \times 4$ spin generators $\sigma^{ab}$ decompose as
$\sigma^{ab} = \sigma^{ab}_L \oplus \sigma^{ab}_R$.
Contracting:
\begin{alignat}{2}
  \eta^i_{ab}\,\sigma^{ab}_L &= 0 &\qquad &(i = 1,2,3), \notag\\
  \eta^i_{ab}\,\sigma^{ab}_R &\neq 0 &\qquad &(i = 1,2,3), \notag\\
  \bar\eta^i_{ab}\,\sigma^{ab}_L &\neq 0 &\qquad &(i = 1,2,3), \notag\\
  \bar\eta^i_{ab}\,\sigma^{ab}_R &= 0 &\qquad &(i = 1,2,3).
\end{alignat}
This is verified by computing, for each $i$, the $2 \times 2$ upper-left
block $P_L (\eta^i_{ab}\sigma^{ab}) P_L$ and the $2 \times 2$ lower-right
block $P_R (\eta^i_{ab}\sigma^{ab}) P_R$, finding the former vanishes and
the latter is a nonzero linear combination of Pauli matrices.  The
computation for $\bar\eta^i$ gives the opposite result.

This establishes the crossed chirality assignment: self-dual Weyl components
$C^+ \sim W^+_{ij}\eta^i\eta^j$ couple to $\Omega^R$ (right-handed spinors),
and anti-self-dual components $C^- \sim W^-_{ij}\bar\eta^i\bar\eta^j$
couple to $\Omega^L$ (left-handed spinors).


% ============================================================
\section{Molien series derivation}
\label{app:molien}
% ============================================================

We derive the parity-even Molien series~\eqref{eq:molien} for
completeness.

The Weyl tensor in $d = 4$ decomposes under $\mathrm{SO}(4)$ as:
\begin{equation}
  W = W^+ \oplus W^- \in \mathrm{Sym}_0^2(\mathbb{R}^3)
  \oplus \mathrm{Sym}_0^2(\mathbb{R}^3),
\end{equation}
where each $\mathrm{Sym}_0^2(\mathbb{R}^3)$ is the space of traceless
symmetric $3 \times 3$ matrices (5-dimensional).  The invariant ring of
a single $\mathrm{Sym}_0^2$ factor under $\mathrm{SO}(3)$ has the
Molien series~\cite{Fulling:1992vm}:
\begin{equation}
  M(t) = \frac{1 + t^6}{(1 - t^2)(1 - t^3)(1 - t^4)}.
\end{equation}
For two copies with parity $W^+ \leftrightarrow W^-$, the parity-even
(``scalar'') invariants have the Molien series:
\begin{equation}
  P(t) = \frac{M(t)^2 + M(t^2)}{2}.
\end{equation}
The first term counts all $\mathrm{SO}(3) \times \mathrm{SO}(3)$ invariants
(product of individual Molien series); the second corrects for the
$\mathbb{Z}_2$ quotient by parity.  Expanding:
\begin{center}
\begin{tabular}{l|ccccccc}
\toprule
Degree in $W$ & 0 & 2 & 3 & 4 & 5 & 6 & 8 \\
\midrule
$P(\text{deg})$ & 1 & 1 & 0 & 3 & 0 & 5 & 13 \\
\bottomrule
\end{tabular}
\end{center}
At degree~4 (quartic Weyl): $P_4 = 3$ parity-even invariants, confirming
the count~\eqref{eq:three-invariants}.

In the loop-order correspondence (degree = $2L + 2$ for an $L$-loop
on-shell counterterm in Ricci-flat gravity), $L = 3$ corresponds to
degree~4, giving 3 invariants versus 1 spectral parameter.  After
Cayley--Hamilton ($3 \to 2$) and chirality ($2 \to 1$), the system becomes
determined.


\bibliographystyle{unsrt}
\bibliography{../references/sct_chirality}


\end{document}
